\documentclass[../../../thesis.tex]{subfiles}
\begin{document}
\subsection{Bibliografické záznamy}
V zozname použitej literatúry v~závere práce sa nachádzajú
podrobné záznamy použitých zdrojov.
Sú to najmä mená autorov, názov článku alebo číslo kapitoly
knihy, názov časopisu alebo knihy, vydavateľ, rok vydania,
strana, na ktorej sa nachádza citovaný článok a podobne.
V~prípade webových stránok je potrebné uviesť internetovú adresu
a~dátum, kedy sme informáciu zo stánky čerpali.
Dôležité je, aby boli informácie na základe týchto detailov
ľahko a jednoznačne dohľadateľné.
Norma STN ISO 690: 2022 tiež definuje formu takéhoto zoznamu~\cite{iso690}.

V~tejto publikácii formátujeme použitú literatúru
podľa štýlu iso-numeric \cite{Hoftich2022iso690} balíka \texttt{biblatex},
\footnote{\href{https://ctan.org/pkg/biblatex}{\texttt{ctan.org/pkg/biblatex}}}
ktorý do veľkej miery rešpektuje doteraz zaužívané zvyklosti
a~je navrhnutý tak,
aby spĺňal odporúčania normy aj v~slovenskom jazyku.

\subsubsection{Príklad záznamu v databázovom súbore .bib}\label{sec:citExample}
Súbor \texttt{bibliography.bib} is textový súbor, ktorý musí mať predpísanú štruktúru.
Môžeme ho vytvoriť ručne alebo použiť niektorý zo systémov na správu publikácií a~citácií ako sú JabRef, Mendeley, Zotero a podobne.

Nasledujúci bibliografický záznam
\begin{trivlist}
  \item STEINEROVÁ, J. Princípy formovania vzdelania
  v~informačnej vede. \textit{Pedagogická revue.} 2000, \textbf{2}(3), 8--16.
\end{trivlist}
vyzerá v zdrojovom súbore \texttt{bibliography.bib} takto:
\begin{verbatim}
@article{Steinerova2000Principy,
  author  = {J. Steinerová},
  journal = {Pedagogická revue},
  title   = {Princípy formovania vzdelania v informačnej vede},
  year    = {2000},
  number  = {3},
  pages   = {8--16},
  volume  = {2},
}
\end{verbatim}

Každý záznam začína symbolom \verb|@| a nasleduje typ publikácie definovaných v dokumentácii k balíčku Bib\LaTeX.
Okrem článku (\verb|article|) to môže byť aj kniha (\verb|book|),
príspevok v zborníku konferencie (\verb|inproceedings|),
webová stránka (\verb|online|),
správa (\verb|techreport|),
všeobecný záznam (\verb|misc|) a mnohé iné.

Prvý povinný údaj je jedinečný identifikátor, ktorý je ľubovoľný.
Odporúča sa však aby obsahoval priezvisko prvého autora, rok vydania a prvé slovo názvu.

Jednotlivé položky databázového záznamu sú oddelené čiarkou,
za posledným poľom už čiarka nie je.

Každý typ záznamu má iné požiadavky na polia.
Niektoré základné vlastnosti opíšeme v nasledujúcom texte.
Viac informácií možno získať v príslušnej dokumentácii~\cite{Hoftich2022iso690}.

\subsubsection{Prvky bibliografického záznamu}
Zoznam použitej literatúry obsahuje informácie o jednotlivých zdrojoch. Je potrebné mať na pamäti, aby záznamy boli jasné, stručné a jednoznačne identifikovateľné. Bežne zorientovaný čitateľ práce, ktorý sa pohybuje v okruhu tém práce by mal byť schopný identifikovať jednotlivé diela na prvé pozretie. Prípadný záujemca o detailné štúdium problematiky by si mal vedieť na základe záznamu vyhľadať citovanú publikáciu buď na internete alebo v knižnici. Týmito pravidlami sa riadime pri vytváraní zoznamu použitej literatúry.

Uvedieme niekoľko jednoduchých a~najčastejších príkladov bežného
spôsobu citovania.
Snaha je, aby sme sa prirodzene naučili formy uvádzania
bibliografických záznamov v~našej oblasti vedy a techniky
a~dokázali vytvoriť správnu databázu zdrojov v systéme Bib\LaTeX.

\subsubsection*{Mená tvorcov}
V zozname literatúry majú formu: PRIEZVISKO, Meno alebo PRIEZVISKO, M. a~navzájom sú oddelené bodkočiarkou.
Za zoznamom autorov nasleduje bodka.

\subsubsection*{\normalsize Pole v .bib súbore}
\begin{verbatim}
author = {Meno1 Priezvisko1 and Meno2 Priezvisko2 and Meno3 Prizvisko3}
\end{verbatim}
Autorov zapisujeme ako \verb|Meno Priezvisko| a~oddeľujeme ich kľúčovým slovom and.
Môžeme tiež použiť obrátený tvar \verb|Priezvisko, Meno and|\dots s~čiarkou medzi priezviskom a~menom.

\subsubsection*{Názov}
Názov článku píšeme normálnym písmom.
Názov nosného informačného zdroja (kniha, časopis, zborník) píšeme kurzívou.

\subsubsection*{\normalsize Polia v .bib súbore}
\begin{verbatim}
title = {}
booktitle = {}
\end{verbatim}

\subsubsection*{Dátum}
Pri každej publikácii musí byť uvedený aspoň rok vydania
vo formáte YYYY, t. j. 2024.
V prípade online zdrojov treba uvádzať aj presný dátum,
kedy sme zdroj použili.
Vkladáme ho do hranatých zátvorkiek s~kľúčovým slovom cit.
Odporúčame použiť formát [cit. YYYY-MM-DD].
Ak sme teda informáciu z webovej stránky získali dňa
14. mája 2023,
napíšeme to takto: [cit. 2023-05-14].
Ide o medzinárodne akceptovaný a~zrozumiteľný tvar zápisu dátumu.

\subsubsection*{\normalsize Polia v .bib súbore}
\begin{verbatim}
year = {2024}
date = {2024-12-31}
\end{verbatim}

\subsubsection*{Dostupnosť}
Myslí sa tým dostupnosť na internete najmä prostredníctvom webu.
Uvádzame buď úplnú webovú adresu alebo DOI číslo.
Nepíšeme oba údaje, preferujeme DOI.
Tento údaj sa nachádza zväčša na konci záznamu za kľúčovým výrazom Dostupné na.

\subsubsection*{\normalsize Polia v .bib súbore}
\begin{verbatim}
url = {}
doi = {}
eprint = {}
\end{verbatim}
Ak uvedieme viacero polí, Bib\LaTeX\ zväčša vyberie iba jedno
podľa preferencií špecifikovaných v~definičnom súbore.

\subsubsection*{Ročník alebo zväzok a číslo časopisu}
Vedecké periodiká vychádzajú v tzv. zväzkoch (angl. \foreignlanguage{english}{\it volume}).
Zväzky sa môžu deliť na čísla
(angl. \foreignlanguage{english}{\it issue} alebo
\foreignlanguage{english}{\it number}).
Zväzok a číslo zapisujeme buď pomocou skratiek vol. a~iss.
(prípadne no. podľa skutočného členenia jednotlivých vydaní konkrétneho časopisu),
alebo preferujeme zaužívanú skrátenú formu,
pri ktorej nepoužijeme slovné skratky,
ale píšeme len čísla,
pričom zväzok vytlačíme tučným rezom (bold)
tesne nasledovaný číslom časopisu normálnym písmom
v~okrúhlych zátvorkách.
Napríklad zväzok 15, číslo 8 môžeme zapísať
ako vol. 15, iss. 8 alebo \textbf{15}(8).

\subsubsection*{\normalsize Polia v .bib súbore}
\begin{verbatim}
volume = {}
number = {}
\end{verbatim}
\end{document}